\documentclass[11pt]{article}
\usepackage[a4paper,margin=27mm]{geometry}
\usepackage[T1]{fontenc}
\usepackage{lmodern,microtype,amsmath,amssymb,amsthm,enumitem}
\usepackage[hidelinks]{hyperref}
\newtheorem{theorem}{Theorem}
\newtheorem{lemma}[theorem]{Lemma}
\theoremstyle{remark}\newtheorem{remark}[theorem]{Remark}
\newcommand{\U}{\mathcal U}
\newcommand{\W}{\mathcal W}
\newcommand{\debt}{\mathrm{debt}}
\newcommand{\cred}{\mathrm{credit}}
\title{Independent Rankings Without Additive Structure}
\author{}\date{11 October 2026}
\begin{document}\maketitle

\noindent
\emph{Fineness Spreads Along Edges} (FSE) asks whether a complete ranking
can satisfy guidance improvement (G) and independence (IB), yet have
neither a nonconstant ordered-group primary nor finitely additive local
extensions off arbitrarily small sets (AExt). The answer is yes. In fact,
the example below has no strictly positive two-sided finitely additive
extension on any non-null region. It also satisfies FSE's conclusions
(LH) and (DH).

This settles that auxiliary question, not the main question. We construct
a ranking at a single world. We do not supply a competitive coherent
selection satisfying (S). In particular this is not a counterexample to
the conjecture that (G), (IB) and (S) force the vacuous model to be
dominated. The arguments here are author-checked, not independently
refereed.

\section{A finite starting point}

We use the following classical existence result. There is a finite set
$J$ and a linear order $\le_0$ on its subsets such that
\begin{enumerate}[label=(\roman*),leftmargin=*]
\item $A\subsetneq B$ implies $A<_0B$;
\item $A\le_0B$ if and only if $A\cup C\le_0B\cup C$, whenever
$C\cap(A\cup B)=\varnothing$;
\item no nonzero vector $(w_j)_{j\in J}$ of nonnegative real numbers
satisfies
\begin{equation}
 A\le_0B\quad\Longrightarrow\quad
 \sum_{j\in A}w_j\le\sum_{j\in B}w_j
 \quad\hbox{for all }A,B\subseteq J.                 \label{eq:noalmost}
\end{equation}
\end{enumerate}
This is a finite comparative-probability order that is not even
\emph{almost representable}. Such examples exist with $|J|=6$: see
Kraft--Pratt--Seidenberg \cite{KPS}, and the definition and six-element
results in Conder--Slinko \cite[Sections 2 and 5]{CS}. Here ``almost''
allows zero weights and ties; merely failing strict representation
would not be enough for our argument.

Property (ii) also allows replacement of a common part. If
$A\cap I=B\cap I$, replacing this common subset of $I$ by any other
subset of $I$ preserves their comparison: first remove the old common
part, then add the new one.

\section{The ranking}

Fix a world $i$. Let $H$ represent binary problems, with the atomless
finite Gaussian measure $\mu$ used in the other notes. Let $H'\subseteq H$
be the mixed problems. All sets and models are identified modulo null
sets. A model $M$ has measurable guidance rank
$r_M:H\to\{0,1,2,3\}$ at $i$, with ranks respectively good, silent, bad,
and contradictory. Smaller is better. Every such rank assignment gives
a model, by prescribing its advice separately on each binary problem.

Partition $H'$ into non-null sets $(E_j)_{j\in J}$. Choose an ultrafilter
$\U_j$ of the measure algebra containing $E_j$. An ultrafilter chooses
exactly one cell of a finite measurable partition. Thus $r_{\U}(M)$,
the unique rank whose level set belongs to $\U$, is well defined.

Well-order all ultrafilters of the measure algebra as
$(\W_\alpha)_{\alpha<\kappa}$. The use of the axiom of choice here is
explicit. These ultrafilters separate distinct elements of the measure
algebra. For $l=3,2,1$, put
\[
 S_l(M)=\{j\in J:r_{\U_j}(M)\ge l\}.
\]
Compare models by the lexicographic order of
\begin{equation}
 \bigl(S_3(M),S_2(M),S_1(M);
                (r_{\W_\alpha}(M))_{\alpha<\kappa}\bigr),       \label{eq:order}
\end{equation}
using $\le_0$ on each of the first three entries and the usual numerical
order on the remaining entries. Thus the first entry at which two
models differ decides their comparison. A least differing entry exists
because the index set is well ordered. Write the resulting ranking as
$\preceq$.

\begin{theorem}\label{thm:GIB}
The ranking \eqref{eq:order} is complete and transitive, respects null
equivalence, and satisfies strict (G) and (IB).
\end{theorem}
\begin{proof}
Lexicographic order gives completeness and transitivity. If
$r_X\le r_Y$ almost everywhere, each $S_l(X)$ is contained in $S_l(Y)$,
and every ultrafilter rank of $X$ is at most that of $Y$. If the
improvement has positive measure, some ultrafilter contains a non-null
cell on which $r_X<r_Y$. Hence at least one entry is strictly better,
and the first differing entry is strictly better. This proves (G).

For (IB), suppose two models have common advice on $B$, and replace
that advice by any other common advice there. Each $\U_j$ either
contains $B$, in which case both ranks at $\U_j$ change together, or
contains its complement, in which case neither rank changes. The sets
$S_l(X),S_l(Y)$ therefore undergo replacement of a common part.
Property (ii) preserves their comparison, including equality, for each
$l$. At every tail entry $\W_\alpha$, an equal pair remains equal and
an unequal pair remains unchanged. The first differing entry and its
direction are consequently preserved. All these statements concern
the measure algebra, so null modifications have no effect.
\end{proof}

\section{No nonconstant additive primary}

A real additive primary on regular models has the form
\begin{equation}
 L(M)=\nu_1(H'\cap\{r_M\ge1\})+
       \nu_2(H'\cap\{r_M\ge2\}),                       \label{eq:primary}
\end{equation}
where $\nu_1,\nu_2$ are finite nonnegative finitely additive charges
vanishing on null sets, and
$L(X)<L(Y)\Rightarrow X\prec Y$. Completeness implies the equivalent
weak-extension condition $X\preceq Y\Rightarrow L(X)\le L(Y)$.

\begin{theorem}\label{thm:noprimary}
Every primary of the form \eqref{eq:primary} is constant. There is also
no nonconstant primary of the same form with values in a linearly
ordered abelian group.
\end{theorem}
\begin{proof}
For $A\subseteq J$, let $E_A=\bigcup_{j\in A}E_j$. Start with the vacuous
model $V$, silent on $H'$, and add debt on $E_A$. Its first coarse entry
is empty and its second is $A$, so these models are ordered exactly by
$\le_0$. The numbers $w_j=\nu_2(E_j)$ must therefore satisfy
\eqref{eq:noalmost}. They are all zero; hence $\nu_2(H')=0$.

Next take the regular model silent on $E_A$ and good on $H'\setminus E_A$.
Its first two coarse entries are empty and its third is $A$. Applying
the same argument to $w_j=\nu_1(E_j)$ gives $\nu_1(H')=0$. Thus $L=0$
on regular models.

For the group-valued claim, the first part of FSE's leading-primary
lemma uses no selection axiom. For clarity, here is the argument.
If a group primary is nonconstant, its maximal value
$M=\nu_1(H')+\nu_2(H')$ is positive. In the convex subgroup generated
by $M$, quotient by the largest convex subgroup not containing $M$.
The quotient is an Archimedean ordered group, so it embeds in
$\mathbb R$, with $M$ mapped to $1$. Composing the charges with this
map gives a nonconstant real primary. Strict inequality in the
quotient implies strict inequality before taking the quotient, so
the primary property is preserved. This contradicts the first part.
\end{proof}

\section{No local extension, even off small sets}

\begin{lemma}\label{lem:small}
Every ultrafilter of this measure algebra contains sets of arbitrarily
small positive $\mu$-measure. If $R$ belongs to the ultrafilter, the
small sets may be chosen inside $R$.
\end{lemma}
\begin{proof}
Partition $R$ into finitely many sets of arbitrarily small measure.
Exactly one belongs to the ultrafilter, and no null set does.
\end{proof}

\begin{theorem}\label{thm:nolocal}
For every non-null $R\subseteq H'$, there is no two-sided weak
extension of the ranking on regular pairs agreeing off $R$ of the form
\[
 L_R(M)=\nu_1(R\cap\{r_M\ge1\})+
        \nu_2(R\cap\{r_M\ge2\}),
\]
with both charges vanishing on null sets and positive on every
non-null subset of $R$. In particular (AExt) fails on every edge region.
\end{theorem}
\begin{proof}
Suppose such charges exist. Decompose $\nu_2$ into its countably
additive and pure parts. As in FSE's pure-set lemma, there is a set
$Z\subseteq R$ of arbitrarily small $\mu$-measure outside which its
pure part vanishes. On $R\setminus Z$, therefore, $\nu_2$ has an
integrable density. It remains strictly positive on every non-null
subset there.

Choose $B_j\in\U_j$ of sufficiently small measure that
\[
 R_0=R\setminus\bigl(Z\cup\textstyle\bigcup_{j\in J}B_j\bigr)
\]
is non-null. No $\U_j$ contains any subset of $R_0$. Hence changing
advice inside $R_0$ leaves all three coarse entries unchanged.
Let $\alpha$ be the first index for which $R_0\in\W_\alpha$.
Every earlier ultrafilter ignores every subset of $R_0$.

By Lemma~\ref{lem:small} and absolute continuity, choose
$B\in\W_\alpha$, $B\subseteq R_0$, with
$0<\nu_2(B)<\nu_2(R_0)/2$. Put $A=R_0\setminus B$.
The two models $X=V+\debt(A)$ and $Y=V+\debt(B)$ first differ at
$\W_\alpha$: $X$ is silent there and $Y$ gives debt. Thus $X\prec Y$.
But
\[
 L_R(X)-L_R(Y)=\nu_2(A)-\nu_2(B)>0,
\]
contradicting weak extension. The two models agree off $R_0\subseteq R$.

For (AExt), the complements of sets whose measures tend to zero inside
a fixed non-null edge region are non-null for all sufficiently large
indices. The result just proved excludes an extension on each of them.
\end{proof}

\section{Even the two non-negligibility conditions hold}

\begin{theorem}\label{thm:LHDH}
For every non-null $R\subseteq H'$ there is a strict improvement of $V$
with non-null debt and credit contained in $R$, and there is a strict
worsening of $V$ with non-null debt and credit contained in $R$.
Consequently the ranking satisfies (LH) and (DH).
\end{theorem}
\begin{proof}
Remove finitely many small $B_j\in\U_j$ to obtain non-null
$R_0\subseteq R$ ignored by every $\U_j$, and let $\W_\alpha$ be its
first tail ultrafilter, as above. Choose $B\in\W_\alpha$, $B\subseteq R_0$,
with $A=R_0\setminus B$ non-null. At the first active entry, credit on
$B$ and debt on $A$ gives rank $0$, while the reverse gives rank $2$;
$V$ has rank $1$. Hence
\[
 V+\debt(A)+\cred(B)\prec V
 \prec V+\debt(B)+\cred(A).
\]
All four sets in these two comparisons are non-null and inside $R$.
\end{proof}

\section{Consequence for the remaining question}

The counterexample answers FSE's Remark~Q positively. Its local failure
is stronger than failure of (AExt), and Theorem~\ref{thm:LHDH} shows
that adding (LH) and (DH) does not repair the proposed alternative.
Thus completeness, (G) and (IB), even with these two consequences of
(S), do not force an additive primary or a local additive extension.

What remains is whether the full selection axiom (S) excludes such
rankings when worlds must share one coherent competitive selection.
The construction supplies no answer to that question. In particular,
it neither proves full coverage (A) possible nor proves it impossible.
It identifies why an argument using (S) itself, beyond (LH) and (DH),
is still needed.

\begin{thebibliography}{9}
\bibitem{KPS} C.~H.~Kraft, J.~W.~Pratt and A.~Seidenberg,
Intuitive probability on finite sets,
\emph{Annals of Mathematical Statistics} 30 (1959), 408--419.
\href{https://doi.org/10.1214/aoms/1177706260}{doi:10.1214/aoms/1177706260}.
\bibitem{CS} M.~Conder and A.~Slinko,
A counterexample to Fishburn's conjecture on finite linear qualitative
probability, \emph{Journal of Mathematical Psychology} 48 (2004), 425--431.
\href{https://doi.org/10.1016/j.jmp.2004.08.004}{doi:10.1016/j.jmp.2004.08.004}.
\end{thebibliography}
\end{document}
